\documentclass{article}
\usepackage[T1]{fontenc}
\usepackage{lmodern}
\usepackage{graphicx}
\usepackage{amsmath,amssymb}
\usepackage[a4paper,margin=2cm]{geometry}
\usepackage[absolute,overlay]{textpos}
\setlength{\TPHorizModule}{1mm}
\setlength{\TPVertModule}{1mm}
\usepackage[indonesian]{babel}
\usepackage{hyperref}
\setlength{\emergencystretch}{2em}
% unit: Halaman 2

\begin{document}

% === Halaman 2 (python/native) ===

Pendahuluan

• Sebagian besar kriptografi kunci-publik ( seperti RSA, ElGamal, Diffie-Hellman) 

menggunakan bilangan bulat yang sangat besar dalam komputasinya.

• Sistem seperti itu memiliki masalah yang signifikan dalam menyimpan, memproses 

kunci dan pesan, dan membutuhkan waktu komputasi yang lama.

• Sebagai alternatif adalah melakukan komputasi berbasis kurva eliptik (elliptic curve). 

• Kriptografi yang menggunakan kurva eliptik dinamakan Elliptic Curve Cryptography 

(ECC). 

• Komputasi dengan kurva eliptik menawarkan keamanan yang sama dengan 

algoritma-algoritma tersebut namun dengan ukuran kunci yang lebih kecil.  

2

Sumber: William Stallings, Cryptography and Network Security

Chapter 10, 5th Edition

\end{document}
